\begingroup
\renewcommand{\arraystretch}{1.45}
\newcommand{\TransscendentArticle}[1]{\par\medskip\Needspace{5\baselineskip}\pdfbookmark[1]{Art. #1}{transscendentes-ii-art-#1}\centerline{[#1.]}\nobreak\smallskip}
\fancyhead[L]{\small ZUR THEORIE DER NEUEN TRANSZENDENTEN. II.}

\Needspace{25\baselineskip}
\TransscendentArticle{2}
Man setze
\[
\begin{gathered}P=1+\frac{x^n\mathbin{.}1-x^{2n+1}}{1+x^n\mathbin{.}1+x^{n+1}}
\\[7pt]\quad{}+\frac{x^{2n}\mathbin{.}1-x^{2n+2}\mathbin{.}1-x^{n+1}}{1+x^n\mathbin{.}1+x^{n+1}\mathbin{.}1+x^{n+2}}
\\[7pt]\quad{}+\frac{x^{3n}\mathbin{.}1-x^{2n+3}\mathbin{.}1-x^{n+1}\mathbin{.}1-x^{n+2}}{1+x^n\mathbin{.}1+x^{n+1}\mathbin{.}1+x^{n+2}\mathbin{.}1+x^{n+3}}
+\text{ usw.}\end{gathered}
\]
\[
\begin{gathered}Q=\frac{x^n}{1+x^n}
+\frac{x^{2n}\mathbin{.}1-x^{n+1}}{1+x^n\mathbin{.}1+x^{n+1}}
\\[7pt]\quad{}+\frac{x^{3n}\mathbin{.}1-x^{n+1}\mathbin{.}1-x^{n+2}}{1+x^n\mathbin{.}1+x^{n+1}\mathbin{.}1+x^{n+2}}
+\text{ usw.}\end{gathered}
\]
\[
R=P-Q
\]


\Needspace{26\baselineskip}
\SourceLeaf{447}
Zunächst bestimme man diese Differenz, indem man das erste, zweite, dritte Glied
und so weiter der Reihe $Q$ vom ersten, zweiten, dritten Glied der Reihe $P$ abzieht. So erhält man
\[
\begin{gathered}R=\frac1{1+x^n}
\\[7pt]\quad{}+\frac{x^n\mathbin{.}1-x^n}{1+x^n\mathbin{.}1+x^{n+1}}
\\[7pt]\quad{}+\frac{x^{2n}\mathbin{.}1-x^n\mathbin{.}1-x^{n+1}}{1+x^n\mathbin{.}1+x^{n+1}\mathbin{.}1+x^{n+2}}
\\[7pt]\quad{}+\frac{x^{3n}\mathbin{.}1-x^n\mathbin{.}1-x^{n+1}\mathbin{.}1-x^{n+2}}{1+x^n\mathbin{.}1+x^{n+1}\mathbin{.}1+x^{n+2}\mathbin{.}1+x^{n+3}}
+\text{usw.}\end{gathered}
\]
Wir bezeichnen diese Reihe mit $\varphi(x,n)$.

\Needspace{22\baselineskip}
Ferner bestimme man jene Differenz $R$, indem man das erste, zweite, dritte
Glied und so weiter der Reihe $Q$ vom zweiten, dritten, vierten Glied und so weiter der Reihe
$P$ abzieht. So erhält man
\[
\begin{gathered}R=1-\frac{x^{2n+1}}{1+x^{n+1}}
\\[7pt]\quad{}-\frac{x^{3n+2}\mathbin{.}1-x^{n+1}}{1+x^{n+1}\mathbin{.}1+x^{n+2}}
\\[7pt]\quad{}-\frac{x^{4n+3}\mathbin{.}1-x^{n+1}\mathbin{.}1-x^{n+2}}{1+x^{n+1}\mathbin{.}1+x^{n+2}\mathbin{.}1+x^{n+3}}
-\text{ usw.}\end{gathered}
\]
\Needspace{7\baselineskip}
oder offenbar
\[
R=1-x^{2n+1}\mathbin{.}\varphi(x,n+1)
\]
\Needspace{7\baselineskip}
folglich
\[
\varphi(x,n)=1-x^{2n+1}\mathbin{.}\varphi(x,n+1)
\]
\Needspace{9\baselineskip}
Dieser Schluss gilt allgemein, solange $n>1$. Unter dieser Einschränkung hat man demnach
\[
\varphi(x,n)=1-x^{2n+1}+x^{4n+4}-x^{6n+9}+x^{8n+16}-\text{ usw.}
\]
\Needspace{9\baselineskip}
Dagegen ist im Fall $n=0$ das letzte Glied von $Q$ nicht als verschwindend zu betrachten. Setzt man es $=T$, so wird der erste Wert von $R$ um
$T$ kleiner sein als der zweite, also
\[
T=1-x\varphi(x,1)-\varphi(x,0)
\]
\Needspace{9\baselineskip}
oder
\[
\begin{gathered}\varphi(x,0)=1-x\varphi(x,1)-T\\[7pt]\quad{}=1-x+x^4-x^9+x^{16}-\ .\ .\ -T\end{gathered}
\]
\Needspace{8\baselineskip}
also, da
\[
\varphi(x,0)=\tfrac12
\]
\Needspace{8\baselineskip}
und
\[
2T=\frac{1-x}{1+x}\mathbin{.}\frac{1-xx}{1+xx}\mathbin{.}\frac{1-x^3}{1+x^3}\mathbin{.}\frac{1-x^4}{1+x^4}\quad\text{usw.}
\]
\Needspace{10\baselineskip}
so erhält man
\SourceLeaf{448}
\[
\begin{gathered}\frac{1-x}{1+x}\,.\,\frac{1-xx}{1+xx}\,.\,\frac{1-x^3}{1+x^3}\,.\,\frac{1-x^4}{1+x^4}\ \text{usw.}
\\[7pt]\quad{}=1-2x+2x^4-2x^9+2x^{16}-\text{usw.}\end{gathered}
\]
\Needspace{17\baselineskip}
Der erste Teil ist hier
\[
\begin{gathered}
=1-x\,.\,1-x^3\,.\,1-x^5\,.\,1-x^7\ \text{usw.}\quad
\\[7pt]\quad{}\frac{1-xx}{1+x}\,.\,\frac{1-x^4}{1+xx}\,.\,\frac{1-x^6}{1+x^3}\ \text{usw.}\\[7pt]
=(1-x)^2(1-xx)(1-x^3)^2(1-x^4)(1-x^5)^2(1-x^6)\ \text{usw.}
\end{gathered}
\]
\Needspace{10\baselineskip}
Bezeichnen wir ferner
\[
1-2x+2x^4-2x^9+2x^{16}\ \text{usw.}\quad\text{mit }Fx
\]
so gilt offenbar
\[
Fx\,.\,F(-x)=(Fxx)^2
\]
\clearpage
\endgroup
